\title{The Era of 1-bit LLMs: All Large Language Models are in 1.58 Bits}

\begin{abstract}
Groundbreaking studies like BitNet~\cite{bitnet} are ushering in the advent of 1-bit Large Language Models (LLMs). In this paper, we propose a new 1-bit LLM variant called \textbf{\text{BitNet b1.58}}, in which all individual model parameters (or weights) are drawn from the ternary set \{-1, 0, 1\}. This approach achieves parity with full-precision Transformer LLMs (using FP16 or BF16) of equivalent size and trained on the same amount of data, attaining comparable results for both perplexity and downstream tasks. Simultaneously, it achieves major gains in cost-efficiency, notably reducing latency, memory usage, throughput demands, and energy expenditure. Importantly, the 1.58-bit LLM introduces a novel scaling law and training methodology for future generations of LLMs that are both high-performing and resource-efficient. It also lays the groundwork for a reimagined computational paradigm, motivating the development of specialized hardware tailored to 1-bit LLMs.
\end{abstract}

\section{The Era of 1-bit LLMs}

Artificial Intelligence has experienced an explosive growth in the scale and performance of Large Language Models (LLMs) over recent years. Although these models excel across an array of natural language processing applications, the ballooning number of parameters presents significant obstacles for practical deployment, as well as concerns regarding economic and ecological costs driven by increased energy use.

One method to mitigate these issues is post-training quantization, which produces low-bitwidth models for inference~\cite{smoothquant, gptq, quip, quip_sharp}. This method lowers the precision of both weights and activations, thus substantially decreasing the storage and computation needed to operate LLMs. Progress has been made from 16-bit precision toward even more compact representations, such as 4-bit models~\cite{gptq, awq}. Despite being broadly adopted in industrial contexts, post-training quantization often falls short of optimal results.

Advancements in 1-bit model designs, notably exemplified by BitNet~\cite{bitnet}, offer a compelling path forward for reducing the operational costs associated with large language models (LLMs) without significant loss in model efficacy. Standard LLMs typically utilize 16-bit floating-point representations (FP16 or BF16), where matrix multiplications constitute the primary computational bottleneck, relying heavily on floating-point arithmetic. The energy and computational expense primarily arises from these floating-point additions and multiplications. Conversely, in BitNet, matrix multiplications depend solely on integer additions, enabling substantial reductions in energy requirements. Since power constraints fundamentally limit computational throughput in many hardware platforms, minimizing energy consumption directly enhances processing speed.

Beyond computational efficiency, another major challenge arises from moving model weights from DRAM to on-chip accelerator memory (such as SRAM) during inference, which can be resource-intensive. Although increasing on-chip SRAM size can boost throughput, such upgrades come at a steep cost relative to DRAM. Compared to their full-precision counterparts, 1-bit LLMs demand significantly less memory capacity and bandwidth, yielding faster and less expensive weight loading from DRAM, and thus enabling speedier and more cost-efficient inference.

In this paper, we propose a noteworthy advancement—a 1-bit LLM variant termed \textbf{\text{BitNet b1.58}}, characterized by parameters that are ternary, drawn from \{-1, 0, 1\}. By introducing the additional value 0 to the original BitNet parameter set, the representation effectively becomes 1.58 bits in binary encoding. \text{BitNet b1.58} preserves the foundational advantages of the original BitNet, most notably a computation scheme that eschews multiplication operations in favor of highly optimized procedures, offering considerable energy conservation and superior efficiency in terms of memory usage, throughput, and latency relative to FP16 LLMs. Moreover, \text{BitNet b1.58} provides two distinct improvements: first, it offers enhanced representational power due to explicit support for feature filtering—enabled by the incorporation of 0 within its weights—which substantially benefits 1-bit LLM performance. Second, empirical results demonstrate that \text{BitNet b1.58} can achieve parity with full-precision (FP16) models with respect to perplexity and downstream task outcomes for models starting at 3B parameters under matched training conditions (including model scale and data quantity).

\section{BitNet b1.58}

\text{BitNet b1.58} builds on the foundational BitNet architecture, which is a variant of the Transformer that substitutes the standard \emph{nn.Linear} layer with a \emph{BitLinear} layer. This model is trained from the ground up, employing 1.58-bit quantized weights and 8-bit activations. Several architectural and procedural refinements differentiate it from the original BitNet, as detailed below.

\paragraph{Quantization Function.} In order to restrict weight values to -1, 0, or +1, we utilize the \emph{absmean} quantization approach. This function begins by normalizing the weight matrix using its mean absolute value, followed by quantizing each element to its nearest value in the set \{-1, 0, +1\}:

\begin{equation}
    \widetilde{W} = \mathrm{RoundClip}(\frac{W}{\gamma + \epsilon}, -1, 1),
\end{equation}

\begin{equation}
    \mathrm{RoundClip}(x, a, b) = \max(a, \min(b, \mathrm{round}(x))),
\end{equation}

\begin{equation}
    \gamma = \frac{1}{nm}\sum_{ij} |W_{ij}|.
\end{equation}

For activation quantization, we apply the same method as BitNet, with one notable distinction: instead of rescaling activations prior to applying non-linear functions into the interval $[0, Q_b]$, activations are uniformly rescaled to $[-Q_b, Q_b]$ on a per-token basis. This adjustment eliminates the need for zero-point quantization. It also simplifies both software implementation and hardware deployment, and in our experiments, this results in only a negligible impact on model performance.

\paragraph{LLaMA-alike Components.}

The architecture introduced in LLaMA~\cite{llama, llama2} has become a standard reference for open-source large language models. To facilitate compatibility and support within the open-source ecosystem, our \text{BitNet b1.58} model adopts core components inspired by LLaMA. This includes using RMSNorm~\cite{rmsnorm}, SwiGLU~\cite{swiglu}, rotary positional embedding~\cite{rope}, and excluding all bias terms. Such design choices allow seamless integration of \text{BitNet b1.58} with widely used open-source platforms (e.g., Huggingface, vLLM~\cite{vllm}, and llama.cpp\footnote{\url{https://github.com/ggerganov/llama.cpp}}) with little modification.

\section{Results}

We conducted a comprehensive comparison between \text{BitNet b1.58} and our FP16 version of \text{LLaMA LLM} across several parameter scales. For equitable benchmarking, all models were pre-trained on the RedPajama dataset~\cite{redpajama} with a corpus size of 100 billion tokens. Evaluation focused on zero-shot performance across multiple language understanding benchmarks, including ARC-Easy~\cite{arc}, ARC-Challenge~\cite{arc}, Hellaswag~\cite{hellaswag}, Winogrande~\cite{winoGrande}, PIQA~\cite{piqa}, OpenbookQA~\cite{openbookqa}, and BoolQ~\cite{boolq}. Additionally, we reported validation perplexity using WikiText2~\cite{wikitext2} and C4~\cite{c4} as reference datasets.

We further assessed GPU memory usage and inference latency for both \text{LLaMA LLM} and \text{BitNet b1.58}. These metrics were obtained leveraging the FasterTransformer\footnote{\url{https://github.com/NVIDIA/FasterTransformer}} library, which is specifically optimized for large language model inference on GPUs. Integration of Ladder's 2-bit kernel~\cite{ladder} supported evaluation of \text{BitNet b1.58}. The principal measurement was the time required to generate each output token, given its primary contribution to overall inference cost.

The results summarized in Table~\ref{tab:ppl} highlight both the perplexity values and resource consumption for \text{BitNet b1.58} and \text{LLaMA LLM}. The findings indicate that, starting at the 3B parameter scale, \text{BitNet b1.58} attains a perplexity comparable to the FP16 \text{LLaMA LLM}, while exhibiting a 2.71-fold increase in generation speed and utilizing 3.55 times less GPU memory. Notably, the 3.9B parameter version of \text{BitNet b1.58} offers a 2.4-fold latency advantage and reduces memory requirements by 3.32 times compared to \text{LLaMA LLM} 3B, while substantially surpassing it in performance.

Detailed zero-shot accuracy outcomes are presented in Table~\ref{tab:zero-shot}. Evaluation protocols were aligned with the methodology in \emph{lm-evaluation-harness}\footnote{\url{https://github.com/EleutherAI/lm-evaluation-harness}}. The data demonstrate a diminishing accuracy disparity between \text{BitNet b1.58} and \text{LLaMA LLM} as model size increases. Crucially, \text{BitNet b1.58} matches the baseline FP16 model's performance from 3B parameters upwards. Consistent with the perplexity trend, \text{BitNet b1.58} at 3.9B outperforms the 3B \text{LLaMA LLM} while maintaining significantly lower requirements for memory and latency. These results suggest that \text{BitNet b1.58} offers a clear Pareto improvement relative to leading LLM approaches.

\paragraph{Memory and Latency}

We expanded our experiments by increasing the model sizes to 7B, 13B, and 70B, and conducted an analysis of their computational costs. Figure~\ref{fig:latency-memory-energy} depicts how both latency and memory usage evolve as models become larger. Notably, the acceleration in inference time becomes more pronounced with increased model capacity; for example, \text{BitNet b1.58} at 70B parameters achieves a 4.1$\times$ speed-up over the \text{LLaMA LLM} baseline. This improvement stems from the fact that the computational demand of \emph{nn.Linear} scales up with the number of parameters. Memory utilization also reflects this trend—since the embedding layer maintains full precision and represents a smaller memory share in larger architectures. Measurements for both latency and memory are based on a 2-bit kernel, indicating that further efficiency gains are possible with additional optimizations.

\paragraph{Energy}

We also assessed the estimated energy expenditure for arithmetic operations in both \text{BitNet b1.58} and \text{LLaMA LLM}. Our evaluation concentrates on the matrix multiplication operations, which constitute the primary energy cost driver for large language models. The breakdown presented in Figure~\ref{fig:energy} reveals that the dominant contribution in \text{BitNet b1.58} arises from INT8 addition, whereas \text{LLaMA LLM} relies on FP16 addition and multiplication. Leveraging the energy estimation framework from~\cite{energycost, pokebnn}, our analysis shows that \text{BitNet b1.58} reduces matrix multiplication energy consumption by a factor of 71.4 on 7nm process technology. Additionally, we provide comprehensive energy cost measurements for processing sequences of 512 tokens. The findings indicate that as the scale of the model increases, the relative energy efficiency of \text{BitNet b1.58} further surpasses that of the FP16 \text{LLaMA LLM}. This is attributable to the growing fraction of operations performed by \emph{nn.Linear} layers in larger models, while the overhead from ancillary components diminishes.

\paragraph{Throughput}

We assess the throughput of \text{BitNet b1.58} in comparison to the \text{LLaMA LLM} with 70B parameters, deploying both across two 80GB A100 GPUs via pipeline parallelism~\cite{gpipe}, which permits \text{LLaMA LLM} 70B to run within available hardware constraints. The batch size was incrementally increased to the maximum supported by GPU memory at a fixed sequence length of 512. According to Table~\ref{tab:throughput}, the 70B variant of \text{BitNet b1.58} can accommodate batch sizes up to 11 times larger than those supported by \text{LLaMA LLM}, leading to an 8.9-fold improvement in throughput.

\textbf{\text{BitNet b1.58} establishes a novel scaling law relating model performance and inference resource requirements.} By referencing Figures~\ref{fig:latency-memory-energy} and \ref{fig:energy}, a size equivalence can be established between 1.58-bit and standard 16-bit models:

\begin{itemize}
    \item A 13B BitNet b1.58 model surpasses a 3B FP16 LLM regarding latency, memory footprint, and energy efficiency.
    \item The 30B BitNet b1.58 model is more effective than a 7B FP16 LLM when comparing latency, memory consumption, and energy usage.
    \item Similarly, the 70B BitNet b1.58 model exceeds the efficiency of a 13B FP16 LLM with respect to latency, memory utilization, and energy requirements.
\end{itemize}

\paragraph{Training with 2T Tokens}

Token count in training remains a fundamental variable in scaling LLM performance. To evaluate the scalability of \text{BitNet b1.58} in terms of training tokens, we trained this model with 2T tokens, employing the StableLM-3B~\cite{StableLM-3B-4E1T} data curation method, mirroring the leading open-source 3B model. Evaluation was carried out on a composite benchmark comprising Winogrande~\cite{winoGrande}, PIQA~\cite{piqa}, SciQ~\cite{sciq}, LAMBADA~\cite{lambada}, and ARC-easy~\cite{arc}. We present zero-shot accuracy metrics in Table~\ref{tab:2t}, computing an average across accuracy and normalized accuracy scores when both are available. Scores for StableLM 3b at 2T tokens are adopted directly from its report. The results indicate that \text{BitNet b1.58} consistently achieves higher performance across all benchmark tasks, highlighting robust generalization in 1.58-bit LLMs.

\section{Discussion and Future Work}

\textbf{1-bit Mixture-of-Experts (MoE) LLMs}

Mixture-of-Experts (MoE) architectures have emerged as an efficient solution for scaling LLMs by reducing computational FLOPs. Nevertheless, their practical adoption is often hindered by significant memory usage and the costs associated with inter-chip communication. Employing 1.58-bit LLMs offers promising remedies to these obstacles. The decrease in memory requirements can lower the number of hardware devices necessary for deploying MoE models. Additionally, it mitigates the bandwidth demands for transferring activations between chips. In ideal scenarios, consolidating entire models onto a single device would eliminate communication overhead altogether.

\textbf{Native Support of Long Sequence in LLMs}

With the widespread deployment of LLMs, processing lengthy sequences has become increasingly vital. The predominant bottleneck in long sequence inference is the substantial memory requirement for KV caching. \text{BitNet b1.58} provides a notable improvement toward natively accommodating longer contexts by compressing activations from 16 bits down to 8 bits—effectively allowing for a doubling of sequence length within existing hardware constraints. Furthermore, lossless compression to as little as 4 bits—or lower in 1.58-bit LLMs—could enable even longer sequences, a prospect we plan to investigate further.

\textbf{LLMs on Edge and Mobile}

Implementing 1.58-bit LLMs can bring significant advancements to the deployment of language models on resource-constrained edge and mobile platforms. Limitations in memory and computational throughput typically curtail the practicality of running LLMs on such devices. Thanks to the minimal memory and reduced energy requirements of 1.58-bit LLMs, it becomes feasible to implement sophisticated language models across a broad array of new applications on these platforms. Such advances can extend and enhance the abilities of edge and mobile hardware, making applications that were previously unviable now attainable. Furthermore, since CPUs constitute the dominant processing hardware in these environments, the efficiency of \text{BitNet b1.58} on CPUs stands to further expand their reach and utility.

\textbf{New Hardware for 1-bit LLMs}

Recent innovations exemplified by Groq\footnote{\url{https://groq.com/}} have illustrated the benefits of tailoring hardware (such as LPUs) for the needs of LLMs. Building on this trajectory, we advocate for the development of novel hardware and system architectures explicitly customized for the operational characteristics of 1-bit LLMs, in light of the new computation schemes introduced with BitNet~\cite{bitnet}.